\documentclass[a4paper]{article}
% Shared preamble for the English probability and statistics documents.
% Based on the original course preamble. Compile with pdfLaTeX.
\usepackage[top=1.5cm,bottom=2cm,left=2cm,right=2cm]{geometry}
\usepackage[T1]{fontenc}
\usepackage{lmodern}
\usepackage[utf8]{inputenc}
\usepackage[english]{babel}

\usepackage{amsmath}
\usepackage{amssymb}
\usepackage{amsthm}
\usepackage{mathrsfs}
\usepackage{mathtools}
\usepackage{etoolbox}
\usepackage{tikz}
\usetikzlibrary{shapes.multipart,arrows.meta}
\usepackage{xcolor}
\usepackage{enumerate}
\usepackage[shortlabels]{enumitem}
\usepackage{hyperref}
\usepackage{version}
\usepackage{bbold}
\usepackage{fontawesome5}
\usepackage{fancyhdr}
\lfoot{}
\usepackage{natbib}
\usepackage{graphicx}
\usepackage{pdfpages}
\usepackage{authblk}
    
\usepackage{blindtext}
\usepackage{multicol}
\usepackage{collcell}
\usepackage{wrapfig}
\usepackage{blkarray}
\usepackage{cancel}
\usepackage{appendix}
\usepackage{float}
\makeatletter
\input{tex/assets/slashbox.sty}
\makeatother
\hypersetup{unicode,colorlinks,citecolor=red,filecolor=black,
  linkcolor=red,urlcolor=cyan}
\setlength{\columnsep}{1cm}
\pagenumbering{gobble}

\def\numtext#1{%
  \ifcase#1\or one\or two\or three\or four\or five\or six\or seven\or eight\or nine\or ten\or
  eleven\or twelve\or thirteen\or fourteen\or fifteen\or sixteen\or seventeen\or eighteen\or nineteen\or twenty\or Lots\fi}

\newcommand{\Byte}[1]{%
  \def\nodecontents{}%
  \foreach[count=\l] \k in {#1} {\xappto{\nodecontents}{\noexpand\nodepart{\numtext{\l}}\k}}%
  \node[rectangle split,rectangle split parts=8,rectangle split horizontal,draw] at (2,2){\nodecontents};}

\newcommand{\D}{{\mathrm{d}}}
\newcommand{\cE}{{\mathcal{E}}}
\newcommand{\cG}{{\mathcal{G}}}
\newcommand{\cX}{{\mathcal{X}}}
\newcommand{\cY}{{\mathcal{Y}}}
\newcommand{\cB}{{\mathcal{B}}}
\newcommand{\cF}{{\mathcal{F}}}
\newcommand{\cO}{{\mathcal{O}}}
\newcommand{\cS}{{\mathcal{S}}}
\newcommand{\cT}{{\mathcal{T}}}
\newcommand{\given}{\, \big| \,}
\newcommand{\cL}{{\mathscr{L}}}
\newcommand{\cD}{{\mathcal{D}}}
\newcommand{\cA}{{\mathscr{A}}}
\newcommand{\cP}{{\mathscr{P}}}
\newcommand{\cR}{{\mathscr{R}}}
\newcommand{\A}{{\mathcal A}}
\newcommand{\one}{\mathbb{1}}
\newcommand{\mathbbm}[1]{\mathbb{#1}}

\providecommand{\C}{}
\renewcommand{\C}{\mathbb C}
\newcommand{\R}{\mathbb R}
\newcommand{\Q}{\mathbb Q}
\newcommand{\N}{\mathbb N}
\newcommand{\T}{\mathbb T}
\newcommand{\F}{\mathbb F}
\renewcommand{\P}{\mathbb P}
\newcommand{\Z}{\mathbb Z}
\newcommand{\E}{\mathbb {E}}
\newcommand{\Pbb}[1]{\P\left( #1\right)}
\newcommand{\Ebb}[1]{\E\left[ #1\right]}

\newcommand{\lb}{\left (}
\newcommand{\rb}{\right )}
\newcommand{\lbb}{\left [}
\newcommand{\rbb}{\right ]}
\newcommand{\labs}{\left |}
\newcommand{\rabs}{\right |}
\newcommand{\lbrb}[1]{\lb #1 \rb}
\newcommand{\lbbrbb}[1]{\lbb#1\rbb}
\newcommand{\labsrabs}[1]{\labs#1\rabs}
\newcommand{\lbbrb}[1]{\lbb#1\rb}
\newcommand{\lbrbb}[1]{\lb#1\rbb}
\newcommand{\langlerangle}[1]{\langle#1\rangle}
\newcommand{\lbcurly}{\{}
\newcommand{\rbcurly}{\}}
\newcommand{\curly}[1]{\lbcurly#1\rbcurly}
\newcommand{\abs}{\labsrabs}
\newcommand{\ind}[1]{1_{#1}}
\newcommand{\indep}{\perp \!\!\! \perp}
\renewcommand{\d}{\, \mathrm{d}}

\newtheorem{theorem}{Theorem}
\newtheorem{problemCnt}{Lemma}
\theoremstyle{definition}
\newtheorem{problem}[problemCnt]{Problem}
\newtheorem{problembook}[problemCnt]{\faBook \,  Problem}
\newtheorem{problempencil}[problemCnt]{\faPencil* \,  Problem}
\newtheorem{problemstar}[problemCnt]{\faStar\,  Problem}
\newtheorem{problemstarpencil}[problemCnt]{\faStar \, \faPencil* \,   Problem}
\newtheorem{probleminteresting}[problemCnt]{\faSmile[regular] \,   Problem}
\newtheorem{solution}{Solution}

\newcommand{\fwcell}[1]{\makebox[\arraystretch\normalbaselineskip]{$#1$}}
\newenvironment{enumeratei}{\begin{enumerate}[(i)]}{\end{enumerate}}
\newenvironment{enumeratea}{\begin{enumerate}[(a)]}{\end{enumerate}}
\newenvironment{enumeraten}{\begin{enumerate}[\upshape 1.]}{\end{enumerate}}

% Compatibility helpers needed by a few files in the combined archive.
\AtBeginDocument{\noindent{\footnotesize English translation prepared with ChatGPT (OpenAI, GPT-6), 26 October 2026, for translation of the Bulgarian text into English and LaTeX formatting.
The Bulgarian originals remain the reference; translation may contain errors.}\par\smallskip}
\newcommand{\MissingGraphic}[1]{%
  \begin{center}\fbox{\parbox{0.82\linewidth}{\centering Missing source image: \texttt{#1}}}\end{center}}



\setlength{\emergencystretch}{3em}
\pagenumbering{arabic}
\begin{document}
\begin{center}\Large Probability and Statistics, Computer Science\par\large Regular exam, 30 June 2022: solutions\end{center}
\noindent\textit{English transcription of the supplied handwritten workings. Numerical conclusions and notation are retained as written.}
\section*{Problem 1}
The completed magic square is
\[\begin{array}{|c|c|c|}\hline2&7&6\\\hline9&5&1\\\hline4&3&8\\\hline\end{array}\]
\begin{enumerate}
\item The expected number of correct entries is $\E[Bin(5,1/10)]=1/2$.
\item The expected number of simulations is $\E[Ge(1/10^5)]=10^5$.
\item Let $X$ be the number of simulations until the first success. With $p=10^{-5}$,
\[X\sim Ge(p),\qquad\E X=10^5=1/p,\qquad DX=(1-p)/p^2.\]
We seek $n$ with $\P(X>n)\le1\%$. A rough Chebyshev bound is
\[\P(X-\E X>\varepsilon)\le\P(|X-\E X|>\varepsilon)\le\frac{DX}{\varepsilon^2}=1\%.\]
Thus take $\varepsilon=10\sqrt{DX}$ and $n=\lceil\E X+10\sqrt{DX}\rceil=1099995$.
\end{enumerate}
For the second simulation procedure:
\begin{enumerate}
\item By indicator variables, the expectation is
\[\E(\one_{\{\text{entry 1 correct}\}}+\cdots+\one_{\{\text{entry 5 correct}\}})=5\E\one_{\{\text{entry 1 correct}\}}=5\cdot\frac15=1.\]
\item $\E[Ge(1/5!)]=5!=120$.
\item As above, choose $n=\lceil\E Y+10\sqrt{DY}\rceil=1315$ for $Y\sim Ge(1/5!)$.
\end{enumerate}
\section*{Problem 2}
Let $Z_i$ be the number on the $i$-th roll. The handwritten calculation is
\begin{align*}
&Cor(Z_1+\cdots+Z_{n-1},Z_2+\cdots+Z_n)\\
&=\frac{\E[(Z_1+\cdots+Z_{n-1})(Z_2+\cdots+Z_n)]-(\E[Z_1+\cdots+Z_{n-1}])^2}{D(Z_1+\cdots+Z_{n-1})}\\
&=\frac{(n-2)\E Z_1^2+[(n-1)^2-(n-2)](\E Z_1)^2-(n-1)^2(\E Z_1)^2}{(n-1)[\E Z_1^2-(\E Z_1)^2]}=\frac{n-2}{n-1}.
\end{align*}
The note says that the result does not depend on how the dice are filled, provided the results are independent and have equal finite expectations and variances.
\section*{Problem 3}
Let $X_i$ be the money collected on day $i$. With $n$ volunteers, $X_i$ is a sum of $10n$ independent variables of the form $100\cdot Ber(1/25)$. Hence
\[\E X_i=40n,\qquad DX_i=n\cdot100^2\cdot\frac1{25}\frac{24}{25}.\]
The notes seek $n$ such that
\begin{align*}
99\%&\le\P(X_1+\cdots+X_{100}\ge10^6)\\
&=\P\left(\frac{X_1+\cdots+X_{100}-4000n}{\sqrt{n\cdot100^2\cdot24/25^2\cdot100}}\ge\frac{10^6-4000n}{\sqrt{n\cdot100^2\cdot24/25^2\cdot100}}\right)\\
&\approx\P\left(N(0,1)\ge\frac{10^6-4000n}{80\sqrt{6n}}\right).
\end{align*}
The next line in the source writes
\[\frac{10^6-4000n}{80\sqrt{6n}}\le\Phi^{-1}(0.01)\approx2.33,\]
followed by $n\ge248.20$, and the boxed conclusion ``$n=249$ is sufficient''. A red annotation says ``250 is enough''. These are transcribed source conclusions; in particular, the displayed sign of the normal quantile is as written in the original.

For the second part, the notes write
\[\P\left(\frac{X_1+\cdots+X_{99}-40\cdot99\cdot249}{\sqrt{n\cdot100^2\cdot24/25^2\cdot99}}>\frac{10^6-40\cdot99\cdot249}{7218.26}\right)\approx\P(N(0,1)\ge4.53)\approx0.\]
\section*{Problem 4}
\begin{enumerate}
\item Normalisation gives
\[1=c\int_0^\infty\int_0^\infty xe^{-(x+y)/2}\,dy\,dx
=2c\int_0^\infty xe^{-x/2}\,dx=8c\int_0^\infty ze^{-z}\,dz=8c,\]
so $c=1/8$. The margin notes $\int xe^{-x}dx=-xe^{-x}-e^{-x}$.
\item
\begin{align*}
\P(X>1,Y>1)&=c\int_1^\infty\int_1^\infty xe^{-(x+y)/2}\,dy\,dx\\
&=2ce^{-1/2}\cdot4\left[-\frac x2e^{-x/2}-e^{-x/2}\right]_1^\infty
=e^{-1/2}\frac32e^{-1/2}=\frac3{2e}.
\end{align*}
\item
\[\E\left[\frac YX\right]=c\int_0^\infty\int_0^\infty\frac yx xe^{-(x+y)/2}\,dx\,dy
=c\int_0^\infty\int_0^\infty ye^{-(x+y)/2}\,dx\,dy=1,\]
by symmetry with the normalisation integral.
\item The marginal density is $f_Y(y)=e^{-y/2}/2$ for $y>0$. Thus
\[f_{X\mid Y}(x,1)=\frac{xe^{-x/2}\cdot\frac18e^{-1/2}}{e^{-1/2}/2}=\frac{xe^{-x/2}}4,\]
and
\[\E[X\mid Y=1]=\int_0^\infty\frac{x^2e^{-x/2}}4\,dx=4.\]
The factorisation $f_{X,Y}(x,y)=c\,xe^{-x/2}e^{-y/2}$ also shows independence.
\end{enumerate}
\end{document}
